\documentclass[11pt]{article}
\usepackage[margin=1in]{geometry}
\usepackage{amsmath,amssymb,amsthm,mathtools}
\usepackage[hidelinks]{hyperref}
\newtheorem{theorem}{Theorem}
\newtheorem{lemma}[theorem]{Lemma}
\newtheorem{corollary}[theorem]{Corollary}
\title{Dyadic unit transfer and next-prime control of cubic binomials}
\author{Internal algebraic proof; no priority or external-review claim}
\date{17 November 2026}
\begin{document}
\maketitle
\begin{theorem}\label{thm:main}
Let $q=2^r$, $r\ge0$. For every integer Gamma shape
$1\le\beta\le2q$, every nonzero
\[
                         f(X)=X^q(A+BX^q)
\]
satisfies Strong Factorial at every starting index. In particular this
holds for ordinary factorials, $\beta=1$, simultaneously for every such
$q$. The coefficients $A,B$ may be arbitrary complex numbers.
\end{theorem}
For $r\ge1$ this has a genuine monomial gap $q\ge2$, outside the
\href{UNIFORM_PADE_PARABOLA_SFC.tex}{completed adjacent-support family}.
The proof below is independent of that parabola argument and of the
constant-plus-power contour or finite-prefix results. It does not prove
all gaps, all lower exponents, all positive Gamma shapes, general FC,
or unrestricted SFC.

\section{A degree-preserving Bessel reduction}
For $n\ge0$ put
\[
 Y_n(z)=\sum_{k=0}^n b_{n,k}z^k,\qquad
 b_{n,k}=\frac{(n+k)!}{2^k k!(n-k)!}
       =\binom{n+k}{2k}(2k-1)!! .                           \tag{1}
\]
The empty odd double factorial is one. Thus every coefficient is integral,
and the leading coefficient $(2n-1)!!$ is odd. Direct coefficient
substitution gives
\[
 Y_0=1,\quad Y_1=1+z,\quad
 Y_n=(2n-1)zY_{n-1}+Y_{n-2}\quad(n\ge2).                  \tag{2}
\]
For interior coefficients the identity reduces to
\[
 (n+k)(n+k-1)-2k(2n-1)-(n-k)(n-k-1)=0;
\]
the constant and extreme coefficients follow directly from (1).
Define also
\[
 D_n=(n+1)Y_n+zY_n'
     =\sum_{k=0}^n\frac{(n+k+1)!}{2^k k!(n-k)!}z^k.        \tag{3}
\]
Its leading coefficient $(2n+1)(2n-1)!!$ is again odd. Differentiating
(2) gives the exact identity
\[
 D_n=(2n-1)zD_{n-1}+D_{n-2}+2Y_n,\quad D_0=1,
                         \quad D_1=2+3z.                 \tag{4}
\]
Consequently both reductions $C_n=\overline Y_n$ and
$C_n=\overline D_n$ in $\mathbb F_2[z]$ have full degree $n$ and satisfy
\[
                        C_n=zC_{n-1}+C_{n-2}.             \tag{5}
\]
Their initial pairs are $(1,1+z)$ and $(1,z)$ respectively.
The Euclidean algorithm in (5) proves
$\gcd(C_n,C_{n-1})=1$ at every index. Their resultants are therefore
nonzero, hence equal to one in $\mathbb F_2$.

\begin{lemma}\label{lem:units}
Let $h\in\{0,1\}$ and let $(u_L)_{L\ge0}$ be any sequence of rational
$2$-adic units. Then the consecutive polynomials
\[
 F_n(t)=\sum_{k=0}^n\binom nk(n+k+h)!u_{n+k}t^k
\]
are coprime over $\mathbb C$ for every $n\ge1$.
\end{lemma}
\begin{proof}
Put $G_n(z)=F_n(z/2)/n!$. Formula (1) shows that its coefficients
belong to $\mathbb Z_{(2)}$, and it reduces to $Y_n$ if $h=0$ and to
$D_n$ if $h=1$. Each rational unit reduces to one. The leading
coefficient stays a unit, so degrees are preserved at every index.
By (5), reduction of $\operatorname{Res}(G_n,G_{n-1})$ is one.
The rational resultant is therefore nonzero; this excludes a common
complex zero. We do not attempt to embed arbitrary complex coefficients
into $\mathbb Q_2$: the conclusion follows from a nonzero rational
resultant of the coefficient polynomials.
\end{proof}
This is a uniform transfer for arbitrary rational unit sequences,
not an assertion that every characteristic-zero coprimality theorem
survives a factorial substitution.

\section{The factorial unit calculation}
Fix $q=2^r$, and for $L\ge0$, $0\le\delta<q$, set
\[
 W_{L,\delta}=\frac{(qL+\delta)!}
                    {2^{(q-1)L}L!\,\delta!}.              \tag{6}
\]
These are odd integers. For $\delta=0$, start with
$W_{L,0}=1$ at $q=1$; doubling $q$ multiplies the previous quotient
by the odd integer
\[
             \frac{(2K)!}{2^K K!}=(2K-1)!!,
                     \qquad K=(q/2)L.
\]
For the remaining offsets,
\[
             W_{L,\delta}=W_{L,0}\binom{qL+\delta}{\delta}.
\]
The binomial factor is odd: for $1\le j\le\delta<q$,
$v_2(qL+j)=v_2(j)$, so its numerator product and $\delta!$ have
identical valuations. This also covers $L=0$, with $W_{0,\delta}=1$.
The unshifted factorial-unit identity is the $p=2$ instance of
Rocks's retained \href{sources/rocks-2015-author-preprint.pdf}{2015 author
draft}, Lemma 3.8; the argument needed here is proved above.
No priority or final-publication scope claim is made.

Write the shape uniquely as
\[
              \beta-1=qh+\delta,\quad h\in\{0,1\},
                                   \quad0\le\delta<q.
\]
For $L\ge0$, (6) gives the exact identity
\[
 (\beta)_{qL}=2^{(q-1)L}\frac{(L+h)!}{h!}
                         \frac{W_{L+h,\delta}}{W_{h,\delta}}. \tag{7}
\]
In particular this is an exponential scaling times $(L+h)!$ times a
rational $2$-adic unit. By Lemma~\ref{lem:units}, it already supplies
the required coprimality. For explicit integer witnesses, define
\[
 M_{n,q,\beta}(t)=\sum_{k=0}^n\binom nk(\beta)_{q(n+k)}t^k,
\qquad
 G_{n,q,\beta}(z)=\frac{h!W_{h,\delta}}{n!2^{(q-1)n}}
                            M_{n,q,\beta}(2^{-q}z).         \tag{8}
\]
Its coefficient of $z^k$ is exactly
\[
               W_{n+k+h,\delta}
                       \frac{(n+k+h)!}{2^k k!(n-k)!}.       \tag{9}
\]
These are integers. They reduce to the coefficients of $Y_n$ or $D_n$,
and their leading coefficient is odd. Thus
\[
             \operatorname{Res}(G_{n,q,\beta},G_{n-1,q,\beta})
                                      \equiv1\pmod2.       \tag{10}
\]
The variable scaling $2^{-q}$ in (8) is the SAME at both moment orders;
only the nonzero overall normalization depends on $n$.

There is also an exact resultant valuation. Since the reduction in (10)
is a unit, the scaling law for resultants in (8) gives
\[
 \boxed{v_2\operatorname{Res}(M_{n,q,\beta},M_{n-1,q,\beta})
 =(3q-2)n(n-1)+(n-1)v_2(n!)+nv_2((n-1)!).}                \tag{11}
\]
Indeed multiplying degree-$n$ and degree-$(n-1)$ polynomials by constants
$c_n,c_{n-1}$ multiplies their resultant by $c_n^{n-1}c_{n-1}^n$;
scaling both variables by $s$ multiplies it by $s^{n(n-1)}$.
Here $v_2(c_j)=-v_2(j!)-(q-1)j$ and $v_2(s)=-q$.
Equivalently, writing $s_2(j)$ for the binary digit sum, (11) is
$3qn(n-1)-(n-1)s_2(n)-ns_2(n-1)$.
The valuation is independent of the particular integer shape in
$1\le\beta\le2q$. This is an exact all-index identity, not an estimate
for a sampled set of resultants.

For $A,B\ne0$, the moment of $X^q(A+BX^q)$ is
$A^n M_{n,q,\beta}(B/A)$. Equation (10) excludes simultaneous zero
moments at $n-1,n$ for every $n\ge1$. If a coefficient vanishes,
$f$ is a nonzero monomial with every positive moment nonzero.
This completes Theorem~\ref{thm:main}.
Multiplying $f$ by any monomial in OTHER independent variables preserves
this two-moment assertion, because those variables contribute nonzero
Gamma factors at each order. Their positive shapes need not be integral.

\section{Limits of the transfer}
The degree-preserving coprime reduction, rather than a characteristic-zero
SFC theorem alone, is the essential input. For example, with $h=2$ the
natural integer polynomials
\[
 T_n(z)=\sum_{k=0}^n\binom{n+k+2}{2}b_{n,k}z^k
\]
have $\overline T_2=\overline T_3=z^2$.
This prevents an automatic extension of this reduction to shapes
$\beta>2q$; it does not refute SFC for those shapes.
Nor does the theorem prove every $X^a(A+BX^2)$. At $(a,n)=(1,3)$,
the ordinary moment polynomials at orders two and three have
\[
 (v_2\text{ coefficients})=(1,4,4),\qquad(1,3,4,7).
\]
Both lower Newton polygons have an edge of slope $3/2$, whose residual
polynomial is $1+Y$. Thus the first residue-field slope test leaves a
common candidate; the actual characteristic-zero polynomials are coprime.
These exact limitations are retained in the checker. More binary digits,
other primes or a different argument could still exclude such candidates.

\section{An exact multiple-root bridge when the lower exponent is one}
For every positive integer gap $e$, put
\[
 A_{n,e}(t)=\frac1{n!}\sum_{k=0}^n\binom nk(n+ek)!t^k,
 \qquad T_{n,e}(t)=\sum_{k=0}^n\binom nk(n)_{ek}t^k,
 \qquad\theta=t\frac{d}{dt}.
\]
For $n\ge1$, direct coefficient comparison gives
\[
 A_{n,e}=T_{n,e}+\frac e n\theta T_{n,e},\qquad
 A_{n-1,e}=T_{n,e}-\frac1n\theta T_{n,e}.                  \tag{12}
\]
The second identity uses both $\binom{n-1}k=(1-k/n)\binom nk$ and
the fact that its Gamma shape is $n$, the same as $T_{n,e}$.
The coefficient matrix of $T,\theta T$ has nonzero determinant.
Since $T_{n,e}(0)=1$, a common zero of the two moments is therefore
EXACTLY a multiple root of $T_{n,e}$. Thus
\[
 X(A+BX^e)\text{ satisfies SFC at every start}
 \iff T_{n,e}\text{ is squarefree for every }n\ge1.         \tag{13}
\]
This proves an equivalence, not the squarefreeness assertion on the right.
It does not assume a fixed-shape Appell derivative identity: the shape
of $T_{n,e}$ is $n$ and varies with the index. It supplies the missing
scalar structure for this particular lower exponent, while the general
lower-exponent problem remains separate.

\section{Conditional next-prime controls for the cubic family}
This section concerns a different support, $f=X(A+BX^2)$, and ordinary
factorials only. The separate
\href{CUBIC_ADJOINT_STABILITY_SFC.tex}{adjoint-stability theorem} proves
this family's all-index SFC assertion for every positive Gamma shape.
The conditional arithmetic controls below retain their own narrower scope.
Write
\[
 C_n(t)=\sum_{j=0}^n\binom nj(n+2j)!t^j,
 \qquad J_n=\gcd_{\mathbb Q[t]}(C_n,C_{n-1}),
\]
where $J_n$ is chosen monic.
\begin{lemma}\label{lem:gap}
Suppose $n\ge6$ and $q$ is prime with
\[
                     n<q\le\frac{3(n-1)}2.                \tag{14}
\]
Put $d=q-n$ and $L=\lfloor(d-1)/2\rfloor$. Then every root of $J_n$
is a $q$-adic unit, and
\[
             \deg J_n\le L,\qquad
             \deg J_n\le2\left\lfloor\frac{d-1}{4}\right\rfloor. \tag{15}
\]
The second bound uses the real place. In particular $d\le4$ certifies
coprimality. No assertion that such a prime exists for every small $n$
is part of this conditional lemma.
\end{lemma}
\begin{proof}
Since $q>n$, all binomial coefficients are $q$-adic units. Also
$q^2>3n$, so coefficient $j$ has valuation
\[
                  v_q\!\left(\binom nj(n+2j)!\right)
                           =\left\lfloor\frac{n+2j}{q}\right\rfloor.
\]
Let $L_0=\lfloor(q-n-1)/2\rfloor$ and
$L_1=\lfloor(2q-n-1)/2\rfloor$.
The lower polygon passes through
\[
                  (0,0),\ (L_0,0),\ (L_1,1),\ (n,2),     \tag{16}
\]
with zero-length or collinear edges deleted. Indeed $L_0,L_1$ are the
last indices of heights zero and one; the slopes in (16) are increasing,
except the one stated collinearity below. The same description applies
to $C_{n-1}$ after replacing $n$ by $n-1$.

For $n=2h$, the two positive-rise widths for $C_n$ and $C_{n-1}$ are
respectively
\[
 \left(\frac{q-1}{2},\ 3h-q+1\right),\qquad
 \left(\frac{q+1}{2},\ 3h-q-1\right).                      \tag{17}
\]
If $d\ge3$, all widths are positive. The two first widths differ by one,
the two last widths differ by two, and every last width is smaller than
every first width: the smallest first minus the largest last is
$3(d-1)/2>0$. Thus none of the nonzero slopes coincide.
If $d=1$, the two edges of $C_n$ are collinear and have reduced slope
$1/h$; those of $C_{n-1}$ have slopes $1/(h+1)$ and $1/(h-2)$.
Here $h\ge3$ by (14), so these are again distinct.

For $n=2h+1$ we have $d\ge2$, and the corresponding widths are
\[
 \left(\frac{q+1}{2},\ 3h-q+2\right),\qquad
 \left(\frac{q-1}{2},\ 3h-q+1\right).                      \tag{18}
\]
They are positive. The first widths differ by one, the last by one,
and the smallest first exceeds the largest last by $(3d-2)/2>0$.
Again no nonzero slopes coincide. Each displayed rise is one.
By multiplicativity of Newton polygons, a common factor can therefore
have only horizontal edges. All its roots are units, and its degree
fits within the shorter horizontal width $L_0=L$.
One of $n,n-1$ is even, and its moment polynomial is strictly positive
for real $t$, since $C_j(t)=\int_0^\infty x^j(1+tx^2)^j e^{-x}dx>0$
for real $t$ whenever $j$ is even.
Thus $J_n$ has no real roots and has even degree. This proves (15).
\end{proof}

The remaining common factor also has an exact residual certificate.
For $d\ge1$ define the integer polynomial
\[
 H_d(t)=\sum_{j=0}^{\lfloor(d-1)/2\rfloor}
             (-1)^j\frac{(d-1+j)!}{j!(d-1-2j)!}t^j.        \tag{19}
\]
Because $n=q-d<q$, reduction of $C_n/n!$ gives
\[
            \overline{C_n/n!}=\overline H_d,\qquad
            \overline{C_{n-1}/(n-1)!}=\overline H_{d+1}.    \tag{20}
\]
To see this, all indices beyond $L$ have positive valuation, whereas
$\binom nj\equiv(-1)^j(d)_j/j!$ and
$(n+1)_{2j}\equiv(d-1)!/(d-1-2j)!$ in the horizontal range.

A monic common factor with unit roots has $q$-integral coefficients and
unit constant term. Dividing either normalized moment polynomial by it
preserves integral coefficients. Its reduction retains its degree and
divides both residuals in (20). Consequently
\[
 \deg J_n\le\deg\gcd_{\mathbb F_q[t]}(\overline H_d,
                                           \overline H_{d+1}). \tag{21}
\]
A residual gcd of degree at most one certifies coprimality by parity;
more generally its permitted divisor degrees can be intersected with
the even degrees in (15). In particular $q\nmid\operatorname{Res}(H_d,H_{d+1})$
is sufficient. This is an exact transfer at every supplied prime window,
not a uniform coverage theorem for all $n$ or a claim that every residual
pair is squarefree or coprime.

For this cubic the multiple-root formulation gives a stronger degree
restriction. Under (14), the Newton polygon of $T_{n,2}$ has horizontal
width $\lfloor d/2\rfloor$ and two positive-rise edges of height ONE.
Their widths are
\[
 \left(\frac{q+1}{2},\frac{3n}{2}-q\right)\quad(n\text{ even}),
 \qquad
 \left(\frac{q-1}{2},\frac{3n-2q+1}{2}\right)\quad(n\text{ odd}). \tag{22}
\]
The widths are positive and strictly decreasing under (14).
These formulas follow from the last indices
$\lfloor(q-n)/2\rfloor$ and $\lfloor(2q-n)/2\rfloor$ at heights zero
and one. Every positive-rise edge is primitive. A factor occurring twice
cannot occupy such an edge, because two copies would require twice
its full width. Consequently
\[
 g\in\mathbb Q[t],\quad g^2\mid T_{n,2}\Longrightarrow
 \deg g\le\left\lfloor\frac d4\right\rfloor,
 \qquad \deg g\le2\left\lfloor\frac d8\right\rfloor.        \tag{23}
\]
The second restriction uses the general positive-Gamma real-place fact:
even powers have no real roots; for odd powers the derivative integral
is strictly positive, so every real root is simple. A repeated rational
factor thus has even degree. In particular $d\le7$ certifies the original
cubic SFC pair through (13).

Here the exact horizontal residual is
\[
 K_d(t)=\sum_{k=0}^{\lfloor d/2\rfloor}
       (-1)^k\frac{d(d+k-1)!}{k!(d-2k)!}t^k,
                 \qquad\overline T_{n,2}=\overline K_d.  \tag{24}
\]
Every rational square divisor reduces with its degree preserved into a
square divisor of $\overline K_d$. Moreover direct coefficient comparison
shows
\[
 K_d=\frac{H_d+2H_{d+1}}3,\qquad
 \theta K_d=\frac d3(H_{d+1}-H_d).                         \tag{25}
\]
Under (14) we have $q\ge3d+3$, so $3d$ is invertible modulo $q$.
As $K_d(0)=1$, the common-root condition for $H_d,H_{d+1}$ is exactly
the multiple-root condition for $K_d$ in the residue field. The useful
uniform target is sufficient prime coverage of these residual square
factor degrees, subject to (23), not an assumed coprimality at every
prime. None of (20)--(25) proves that coverage for every index.

The universal proofs are the displayed identities and inequalities,
not bounded parameter scans.
Independent replay, formal verification, external review and novelty
claims remain false. General FC and unrestricted SFC remain open.
\end{document}
